% !TeX root = ../papers/nc-p3.tex
\section{重複二次因子と階数一総次数六の排除}
\label{sec:rank-one-total-five}
\begin{theorem}\label{thm:rank-one-total-five}
四辺有限の標準有限Frobenius窓で $X$ は被約、反対根 $N$ は一般可逆で全二次式を成分に持つとする。
行列式の重複既約因子は存在してもただ一つの一次式である。
階数一軌跡の曲線部分の被約総次数は5以下である。
\end{theorem}
\begin{proof}[重複二次因子]
前節により可能な重複因子は次数1または2で、高々一つである。既約二次式 $q$ が重複すれば
残余四次式 $G$ は $q$ の外の少なくとも六つの印を三重に通る。既約四次は共線な三印まで、
二既約二次の積は二つの頂点までなので $G$ は平面因子を持つ。
残った三次が既約なら、$G$ の三重点はその三次の特異集合に入る。
孤立特異点は四つ以下、特異直線上は三印以下なので不可能である。
従って $G$ は平面を持ち残りも次数2以下である。全印での重複度の和は八次式全体で
$2\cdot15+2\cdot7+15=59$ 以下、必要な60に反する。特に階数二以下の曲面成分が存在すれば唯一の平面である。
\end{proof}
\begin{proof}[例外出力平面に対応する直線]
総次数6なら命題\ref{prop:degree-six-rank-one-component-partitions}により曲線 $C$ の外に一つの印 $p_0$ がある。
その出力を $y_0$ とする。定理\ref{thm:rank-one-curve-degree-bound}の例外平面の議論により、
唯一の階数二曲面 $S$ の核直線族が $\Pi=y_0^\perp$ を覆う。$S$ は今証明したように平面である。
従って $S$ の一般点で行列の像二平面は固定ベクトル $y_0$ を含む。

命題\ref{prop:rank-two-plane-conic-trace}の標準表示は、この固定ベクトルからも直接得られる。
実際 $y_0$ を除いた商での像直線を $b$ とし、商上の三行対称行列を $d b b^t$ と原始に書く。
この行列は恒等零でない。零なら非零成分は第一行・列の四式のみで平面上の二次式の包は六次元にならない。
$\deg b\le1$ で、包が三次元に満たなければ元の行列の像が固定三空間に入り、
定理\ref{thm:opposite-column-divisor}に反する。ゆえに $b$ は入力平面の独立な三座標であり $d$ は定数である。
混合行も $b$ の一次式倍で、平方完成により $N|_S=aa^t+q y_0 y_0^t$、$a$ の商は線型同型となる。
特に $S$ 上の階数一集合は一つの二次曲線の台で、そこに載る $C$ の成分の総次数は高々2である。

一般の $v\in\Pi$ について $Nv=0$ の $S$ 上の解は直線 $J_v=V(a^t v)\cap S$ である。
また $S$ の外に基点曲線はない。もし一般の $v\in\Pi$ でそのような曲線があれば全入射の別の
三次元成分が $\Pi$ を支配して一般繊維次元1を持つ。階数一曲線上の成分は核平面が変わるためこれに該当せず、
階数二の曲面成分は既に $S$ が唯一だからである。一方 $v^\perp$ は $C'$ を六つの異なる点で切る。
$S$ に含まれない各曲線成分と $S$ の有限交点の出力は $y_0$ と異なり、一般のそのような平面は全てを避ける。
従って六つの対応点のうち少なくとも四つは $S$ の外、特に $J_v$ の外にある。
\end{proof}
\begin{proof}[一本の基点直線の外の点数]
四つの独立な二次式の共通基点の曲線台が一本の直線 $J$ だけなら、その外の相異なる基点は三つ以下である。
実際、四点があればうち三点は共線でない（第三の点で別の直線が基点になる）。
四点全部が $J$ を含む同じ平面に入れば、各二次式の平面制限は $J$ の式と一次式の積で、
その一次式は非共線な点を消すから恒等零になる。
そうでなければ四点は $\langle J,p\rangle$ の値で二群以上に分かれる。
各 $p$ に対し別群の $p'$ との平面 $\langle J,p'\rangle$、残り二点を通り $p$ を避ける平面の積を取ると
$J$ と他の三点で消え $p$ で消えない。従って四点は七次元の $H^0(\mathcal I_J(2))$ に四独立条件を課し、
四つの独立な二次式を残せない。これを $Nv$ と $J_v$ に適用して矛盾を得る。
\end{proof}
